\section{Graph comparisons and cycle ranks}
\label{sec:comparisons}

We work over $\mathbb C$.  We will reduce termination to controlling how
often the positive model of a step contains a surface above the base
locus where either morphism fails to be an isomorphism.  The graph
comparison identifies exactly where discrepancies improve, including
in a mixed step.  Ranks of algebraic cycle classes then remove the
steps that lose divisors; after these positive-side surfaces have
disappeared, a second rank comparison removes the corresponding
negative-side surfaces as well.

We use the negativity lemma in the following form: if $q:T\to V$ is
projective birational, $E$ is an $\mathbb R$-Cartier divisor, $-E$ is
$q$-nef, and $q_*E\geq0$, then $E\geq0$; see
\cite[Lemma~3.39]{KM98}.  In particular, an exceptional divisor which
is relatively numerically trivial is zero, by applying the lemma to
both signs.  All varieties in this section are projective.

\subsection{The full bad loci}

Consider a birational diagram
\begin{equation}\label{eq:step-diagram}
 X^-\xrightarrow{\,f^-\,} Z
 \xleftarrow{\,f^+\,}X^+,
\end{equation}
where the varieties are normal fourfolds, the morphisms have connected fibres, $f^+$ is small, and the divisors
$D^\pm=K_{X^\pm}+B^\pm$ are $\mathbb Q$-Cartier. The boundary on $X^+$ is the strict transform of $f^-_*B^-$, and
$-D^-$ and $D^+$ are ample over $Z$. The arguments here do not require the positive pair to be lc in advance.

Let $U\subset Z$ be the largest open over which both morphisms are isomorphisms, and put
\begin{equation}\label{eq:bad-loci}
 \Sigma=Z\setminus U,\qquad
 W^-=(f^-)^{-1}\Sigma,\qquad W^+=(f^+)^{-1}\Sigma,
\end{equation}
with reduced structures. These sets include the inverse image of the bad locus of the other morphism. In particular, $W^+$ can be nonempty when $f^+$ is an isomorphism.

\begin{lemma}\label{lem:graph}
Let $G$ be the normalization of the graph of the birational map in \eqref{eq:step-diagram}, with projections $u:G\to X^-$ and $v:G\to X^+$. Then
\[
 H=u^*D^- -v^*D^+\geq0,
 \qquad
 \operatorname{Supp}H=(G\to Z)^{-1}\Sigma.
\]
The divisor $H$ is exceptional over $X^+$, and $-H$ is ample over $Z$. For every divisorial valuation $\xi$ of the common function field,
\begin{equation}\label{eq:discrepancy-comparison}
 a(\xi;X^+,B^+)-a(\xi;X^-,B^-)=\xi(H)\geq0.
\end{equation}
The inequality is strict exactly when $c_{X^-}(\xi)\subset W^-$, equivalently when $c_{X^+}(\xi)\subset W^+$.

For a nontrivial step,
\begin{equation}\label{eq:graph-dimension}
 \dim W^-+\dim W^+\geq3.
\end{equation}
Always $\dim W^+\leq2$. If $f^-$ is also small, then $\dim W^-\leq2$ and $\dim\Sigma\leq1$.
\end{lemma}

\begin{proof}
Every prime divisor on $X^+$ maps to a prime divisor on $Z$, since $f^+$ is small. Its strict transform on $X^-$ has the same adjoint coefficient: use the compatible canonical divisors and the stated boundary transform. Thus $v_*H=0$. The normalization map from $G$ to its closed graph in $X^-\times_Z X^+$ is finite. The sum of the pullbacks of the relatively ample divisors $-D^-$ and $D^+$ is relatively ample on that product; its finite pullback is $-H$. In particular, $-H$ is $v$-nef, and negativity gives $H\geq0$.

For a proper birational morphism onto a normal variety, its nonisomorphism locus on the target is exactly its positive-dimensional-fibre locus. Indeed, over the zero-dimensional-fibre locus the morphism is finite locally on the target, hence an isomorphism by normality. Consequently each graph fibre over $\Sigma$ has positive dimension: it maps surjectively to each of the two fibres and at least one of those fibres has positive dimension. Graph fibres are connected, by Stein factorization for the proper birational map $G\to Z$.

Every closed point of a connected positive-dimensional projective fibre lies on a curve in that fibre. To see this, a point which lies on no positive-dimensional irreducible component would be an isolated zero-dimensional component, contradicting connectedness; a positive-dimensional projective component contains a curve through each of its points. If a point of the graph fibre were outside $\operatorname{Supp}H$, choose such a curve through it. That curve is not contained in the effective $\mathbb Q$-Cartier divisor $H$, so its intersection with $H$ is nonnegative. This contradicts the relative ampleness of $-H$. Thus the whole inverse image of $\Sigma$ lies in the support. On the common isomorphism open the adjoints agree, so the reverse inclusion holds as well.

Pull the divisor identity back to a smooth common model carrying $\xi$. Compatible canonical divisors give \eqref{eq:discrepancy-comparison}. The order of an effective $\mathbb Q$-Cartier divisor at a valuation is positive exactly when its center lies in the support. Since this support is the full inverse image of $\Sigma$, the two descriptions of strictness follow.

For a nontrivial step the support is nonempty and contains a prime divisor of the fourfold $G$. This divisor has dimension three. Its map into $W^-\times W^+$ is finite onto its image, giving \eqref{eq:graph-dimension}.

The assertion about $W^+$ also follows from $v_*H=0$: a prime divisor in $W^+$ would have a divisorial strict transform on $G$ lying in $\operatorname{Supp}H$, contrary to the zero pushforward. If $f^-$ is small, codimension-one comparison similarly gives $u_*H=0$, so $\dim W^-\leq2$. Finally each component of $\Sigma$ is a positive-dimensional-fibre locus for at least one of the two morphisms. Its full inverse image on that side has dimension at least one larger, so $\dim\Sigma\leq1$.
\end{proof}

\begin{corollary}\label{cor:singularities-preserved}
Under the hypotheses of Lemma~\ref{lem:graph}, lc and klt are preserved. For a small step, an effective terminal pair remains terminal. At a divisor-losing step the discrepancy of every lost prime divisor increases strictly.
\end{corollary}

\begin{proof}
Use \eqref{eq:discrepancy-comparison}. In a small step the same valuations are exceptional on both models. A lost prime divisor is contained in $W^-$.
\end{proof}

\subsection{Ranks of cycle classes}

The use of algebraic cycle classes to count fourfold flips goes back to
\cite{KMM87}; see also \cite[Lemma~10, preprint version]{FujinoCanonical}.
We use the same homological principle for the full bad loci, so that it
also detects divisor loss in a mixed step.

For a projective variety $V$ define
\[
 \rho_r(V)=
 \dim_{\mathbb Q}
 \operatorname{span}\{[A]\in H_{2r}(V,\mathbb Q):
                  A\subset V\text{ irreducible},\ \dim A=r\}.
\]
For an open variety use Borel--Moore homology in the same definition. On a normal projective threefold $S$ we write
$\rho_*(S)=\rho_2(S)$. This is a rank of divisor cycle classes; we do not identify it with a Picard number on a singular variety.

\begin{lemma}\label{lem:cycle-rank}
Let $V$ be projective and $A\subset V$ closed of dimension at most $r$. Restriction to $V\setminus A$ maps the span of $r$-cycle classes onto the corresponding Borel--Moore cycle span. Its kernel is exactly the span of the classes of the $r$-dimensional irreducible components of $A$. In particular,
\begin{enumerate}[label=\textup{(\roman*)}]
\item removing a set of dimension less than $r$ does not change this rank;
\item removing a set with an $r$-dimensional component strictly lowers the rank.
\end{enumerate}
For a sequence of diagrams as in Lemma~\ref{lem:graph}, only finitely many steps lose a prime divisor. On a small tail with no positive-side bad surfaces, only finitely many steps have a negative-side bad surface.
\end{lemma}

\begin{proof}
Every irreducible closed $r$-dimensional subvariety of the open complement is the restriction of its closure in $V$. Hence the restriction of cycle spans is onto. The localization sequence
\[
 H_{2r}(A,\mathbb Q)\longrightarrow H_{2r}(V,\mathbb Q)
 \longrightarrow H^{\mathrm{BM}}_{2r}(V\setminus A,\mathbb Q)
\]
is exact. Since $\dim A\leq r$, its first group is generated by the fundamental classes of its $r$-dimensional components. These images already belong to the cycle span, proving the kernel assertion. If such a component exists, its class in $V$ is nonzero: intersection with the $r$th power of an ample divisor has positive degree.

In a step the complements of $W^-$ and $W^+$ are isomorphic. Apply the first assertion with $r=3$. The positive bad set has dimension at most two, so the positive-side rank equals the rank on this common open. The negative-side rank is at least that rank, and is strictly larger if a divisor is lost. These nonnegative integer ranks cannot strictly decrease infinitely often. This proves the first finiteness assertion, including mixed steps.

After passing to a small tail, both bad sets have dimension at most two. If the positive sets have dimension at most one, repeat the argument with $r=2$. Each negative-side surface gives a strict rank decrease.
\end{proof}

\begin{corollary}\label{cor:surfaces-suffice}
To prove termination of a fourfold program it suffices to show that, on its small tail, positive-side bad surfaces occur only finitely often.
\end{corollary}

\begin{proof}
After the two applications of Lemma~\ref{lem:cycle-rank}, both bad sets have dimension at most one. Their dimensions then contradict \eqref{eq:graph-dimension} at every further nontrivial step.
\end{proof}

\subsection{Identifying a prescribed ample end}

The following elementary comparison will be used for both terminal and dlt bridges. All divisor equalities are equalities of actual $\mathbb Q$-divisors with compatible canonical choices.
It is the usual comparison of nef birational ends by negativity;
see also \cite[Remark~2.7]{BirkarLC}.

\begin{lemma}\label{lem:two-ends}
Let $V_1$ and $V_2$ be normal projective birational models over a normal base $Z$, with relatively nef $\mathbb Q$-Cartier divisors $D_1,D_2$. On a smooth common model $W$, with projections $q_j:W\to V_j$, suppose
\[
 A=q_1^*D_1+F_1=q_2^*D_2+F_2,
\]
where $F_j\geq0$ is exceptional over $V_j$. Then
$q_1^*D_1=q_2^*D_2$. If $D_2$ is ample over $Z$, the birational map $V_1\dashrightarrow V_2$ is a morphism $h$, and $D_1=h^*D_2$.
\end{lemma}

\begin{proof}
Put $E=q_1^*D_1-q_2^*D_2=F_2-F_1$. Its pushforward to $V_1$ is effective. On curves contracted by $q_1$, the divisor $E$ is the negative of a nef divisor, so $-E$ is $q_1$-nef. Negativity gives $E\geq0$. Interchanging the two ends gives $-E\geq0$, hence $E=0$.

If $D_2$ is relatively ample, every curve in a fibre of $q_1$ has zero degree against $q_2^*D_2$, and is therefore contracted by $q_2$. The image under $q_2$ of that projective connected fibre must be a point. Indeed a positive-dimensional projective image would contain a curve, and projective hyperplane sections of its inverse image supply a curve mapping onto it. The graph factorization lemma for a proper morphism with connected fibres onto a normal variety gives $q_2=hq_1$. Equality of pullbacks and birational pushforward yield $D_1=h^*D_2$.
\end{proof}

For a finite negative birational program extracting no prime divisor, repeated use of Lemma~\ref{lem:graph} writes the input adjoint pullback as the final pullback plus an effective final-exceptional divisor. Thus Lemma~\ref{lem:two-ends} applies to finite programs and a prescribed ample model whenever both comparisons have this form.
